\documentclass[10pt]{article}
\usepackage[T1]{fontenc}
\usepackage[utf8]{inputenc}
\usepackage[french]{babel}
\usepackage{amsmath,amssymb,mathrsfs,mathpazo}
\usepackage{geometry}
\geometry{a4paper,margin=2cm}
\usepackage{graphicx,array,tabularx,tabularray,colortbl,multirow,multicol,booktabs,enumitem,pifont,ulem,fancybox,diagbox,cancel}
\usepackage{pstricks,pstricks-add,pst-plot,pst-tree,pst-node,pst-text,pst-func,pst-math,pst-bspline,pst-3dplot,pst-eucl,pst-solides3d}
\let\pscalculate\undefined
\usepackage{pst-bezier}
\usepackage{tikz,pgfplots,tkz-tab}
\usetikzlibrary{arrows,arrows.meta,patterns,calc,positioning,trees,decorations.markings,shapes.geometric,shapes.misc,shapes.arrows,intersections,decorations.pathmorphing,decorations.pathreplacing,backgrounds,shadows,babel}
\pgfplotsset{compat=1.18}
\usepgfplotslibrary{fillbetween}
\usepackage[np]{numprint}
\usepackage{listings,makecell,pas-tableur,twoopt}
\let\barre\undefined
\usepackage[pstricks]{tablvar}
\definecolor{darkgreen}{RGB}{0,100,0}
\providecommand{\red}{\color{red}}
\providecommand{\blue}{\color{blue}}
\providecommand{\green}{\color{green}}
\providecommand{\black}{\color{black}}
\providecommand{\white}{\color{white}}
\providecommand{\gray}{\color{gray}}
\providecommand{\cyan}{\color{cyan}}
\providecommand{\magenta}{\color{magenta}}
\providecommand{\yellow}{\color{yellow}}
\providecommand{\violet}{\color{violet}}
\providecommand{\orange}{\color{orange}}
\providecommand{\euro}{\texteuro}
\providecommand{\EUR}{\texteuro}
\providecommand{\ieme}{\textsuperscript{e}}
\providecommand{\degres}{\ensuremath{^\circ}}
\providecommand{\degree}{\ensuremath{^\circ}}
\providecommand{\vv}[1]{\overrightarrow{#1}}
\providecommand{\Coord}[1]{#1}
\providecommand{\decosix}{}
\providecommand{\starredbullet}{\ensuremath{\bullet}}
\pagestyle{empty}
\setlength{\parindent}{0pt}
\usepackage{adjustbox}\DeclareUnicodeCharacter{2236}{:}\DeclareUnicodeCharacter{2190}{\ensuremath{\leftarrow}}
\providecommand{\R}{\mathbb{R}}
\providecommand{\N}{\mathbb{N}}
\providecommand{\Z}{\mathbb{Z}}
\providecommand{\Q}{\mathbb{Q}}
\providecommand{\D}{\mathbb{D}}
\providecommand{\C}{\mathbb{C}}
\providecommand{\e}{{\mathrm{e}}}
\providecommand{\vect}[1]{\overrightarrow{#1}}
\providecommand{\vectt}[1]{\overrightarrow{\mathrm{#1}}}
\providecommand{\vvt}[1]{\overrightarrow{\mathrm{#1}}}
\providecommand{\np}[2][]{{#2~#1}}
\providecommand{\vv}[1]{{\overrightarrow{#1}}}
\renewcommand{\d}{{\mathrm{d}}}
\providecommand{\degres}{{^\circ}}
\providecommand{\degre}{{^\circ}}
\providecommand{\blue}{\color{blue}}
\providecommand{\red}{\color{red}}
\providecommand{\enskip}{\quad}
\providecommand{\llbracket}{\lbrack\!\lbrack}
\providecommand{\rrbracket}{\rbrack\!\rbrack}
\providecommand{\euro}{{\text{€}}}
\providecommand{\no}{{\text{n°}}}
\providecommand{\interval}[2]{{[#1;#2]}}
\providecommand{\intervFO}[2]{{[#1;#2[}}
\providecommand{\dotfill}{\ldots}

\providecommand{\ds}{\displaystyle}
\providecommand{\textding}[1]{#1}
\providecommand{\Oij}{\(\mathrm{O};\vec{\imath},\vec{\jmath}\)}
\providecommand{\Oijk}{\(\mathrm{O};\vec{\imath},\vec{\jmath},\vec{k}\)}
\providecommand{\Ouv}{\(\mathrm{O};\vec{u},\vec{v}\)}
\providecommand{\intervFF}[2]{\ensuremath{[#1;#2]}}
\renewcommand{\euro}{\ifmmode\text{\texteuro}\else\texteuro\fi}
\DeclareUnicodeCharacter{00B0}{\ifmmode^{\circ}\else\textdegree\fi}
\DeclareUnicodeCharacter{20AC}{\ifmmode\text{\texteuro}\else\texteuro\fi}
% Énoncés : Éducation nationale. Transcription et corrigés : APMEP, Denis Vergès et contributeurs. Découpage et classement : CoopMaths, Sébastien Lozano, Éric Elter et contributeurs. Adaptation de la présentation : Maths Explorer.
% CC BY-SA : https://coopmaths.fr/www/about/
% https://coopmaths.fr/alea/static/bac/2023/tex/bac_2023_03_sujet2_ameriquenord_3_cor.tex
\begin{document}
\noindent\textit{APMEP / CoopMaths. CC BY-SA. Présentation adaptée par Maths Explorer.}\par\medskip
\medskip

Dans l'espace muni d'un repère orthonormé d'unité $1 \mathrm{~cm}$, on considère les points

\[\mathrm{D}(3~;~1~;~5), \qquad 
\mathrm{E}(3~;~-2~;~-1), \qquad 
\mathrm{F}(-1~;~2~;~1), \qquad 
\mathrm{G}(3~;~2~;~-3)\]



\begin{enumerate}
\item 
	\begin{enumerate}
		\item Le vecteur $\overrightarrow{\text{EF}}$ a pour coordonnées
$\begin{pmatrix}
-1-3 \\ 2-(-2) \\ 1-(-1)
\end{pmatrix}
=
\begin{pmatrix}
-4 \\ 4 \\ 2
\end{pmatrix}$
		
		
 Le vecteur $\overrightarrow{\text{FG}}$ a pour coordonnées
$\begin{pmatrix}
3-(-1) \\ 2-2 \\ -3-1
\end{pmatrix}
=
\begin{pmatrix}
4 \\ 0 \\-4
\end{pmatrix}$

		\item 
$-4\times (-1)=4$ et $4\times (-1)=-4\neq 0$ donc les vecteurs $\vectt{EF}$ et $\vectt{FG}$ ne sont pas colinéaires; on en déduit que les points E, F et G ne sont pas alignés.		
		
		
	\end{enumerate}
\item 
	\begin{enumerate}
		\item 
La droite (FG) passe par le point F et a pour vecteur directeur $\vectt{FG}$; elle a donc pour représentation paramétrique:

$\left \{
\begin{array}{l !{=} l l}
x & x_{\text F} + 4t\\
y & y_{\text F} +0t & t\in\R\\
z & z_{\text F} + (-4)t
\end{array}
\right .	$ soit
$\left \{
\begin{array}{l !{=} l l}
x & -1 + 4t\\
y & \phantom{-}2 & t\in\R\\
z & \phantom{-}1-4t
\end{array}
\right.$

		\item On appelle $\mathrm{H}$ le point de coordonnées $(2~;~2~;~-2)$.

$\mathrm{H}$ est le projeté orthogonal de E sur la droite (FG) si H appartient à la droite (FG) et si les vecteurs $\vectt{EH}$ et $\vectt{FG}$ sont orthogonaux.

\begin{list}{\textbullet}{}
\item Le point H appartient à la droite (FG) si on peut trouver une valeur de $t$ telle que:
$\left \{
\begin{array}{l !{=}  l}
x_{\text H} & -1 + 4t\\
y_{\text H} & \phantom{-}2 \\
z_{\text H} & \phantom{-}1-4t
\end{array}
\right .	$
c'est-à-dire
$\left \{
\begin{array}{r !{=}  l}
2 & -1 + 4t\\
2 & \phantom{-}2 \\
-2 & \phantom{-}1-4t
\end{array}
\right .	$

La valeur $t=\dfrac{3}{4}$ convient donc le point H appartient à la droite (FG).

\item Le vecteur $\overrightarrow{\text{EH}}$ a pour coordonnées
$\begin{pmatrix}
2-3 \\ 2-(-2) \\ -2-(-1)
\end{pmatrix}
=
\begin{pmatrix}
-1 \\ 4 \\-1
\end{pmatrix}$
et $\vectt{FG}$ a pour coordonnées
$\begin{pmatrix}
4 \\ 0 \\ -4
\end{pmatrix}$

$\vectt{EH}\cdot\vectt{FG} = (-1)\times 4 + 4\times 0 + (-1)\times (-4)=0$ donc $\vectt{EH}\perp\vectt{FG}$
\end{list}

Le point H est donc le projeté orthogonal du point E sur la droite (FG).

		\item 
D'après les questions précédentes, [EH] est la hauteur du triangle EFG correspondant à la base [FG]; donc l'aire du triangle EFG vaut, en cm$^2$:

$\mathcal{A}=\dfrac{\text{EH}\times \text{FG}}{2}$.		

\begin{list}{\textbullet}{}
\item $\vectt{EH}\;\begin{pmatrix}
-1 \\ 4 \\-1
\end{pmatrix}$
donc
$\text{EH}=\ds\sqrt{(-1)^2+4^2+(-1)^2} = \ds\sqrt{18}=3\sqrt{2}$
\item $\vectt{FG}\;\begin{pmatrix}
4 \\ 0 \\-4
\end{pmatrix}$
donc
$\text{FG}=\ds\sqrt{4^2+0^2+(-4)^2} = \ds\sqrt{32}=4\sqrt{2}$
\end{list}

Donc l'aire du triangle est en cm$^2$: 
$\mathcal{A}= \dfrac{3\sqrt{2}\times 4\sqrt{2}}{2}=12$.

	\end{enumerate}
\item
	\begin{enumerate}
		\item Le vecteur $\vec{n}\begin{pmatrix}2 \\ 1 \\ 2\end{pmatrix}$ est un vecteur normal au plan (EFG) s'il est orthogonal à $\vectt{EF}$ et à $\vectt{FG}$.
		
\begin{list}{\textbullet}{}
\item $\vect{n}\cdot\vectt{EF} = 2\times(-4) + 1\times 4 +2\times 2=0$ donc $\vect{n}\perp\vectt{EF}$
\item $\vect{n}\cdot\vectt{FG} = 2\times 4 + 1\times 0 +2\times (-4)=0$ donc $\vect{n}\perp\vectt{FG}$
\end{list}

Le vecteur $\vec{n}\begin{pmatrix}2 \\ 1 \\ 2\end{pmatrix}$ est donc un vecteur normal au plan (EFG).
		\item
$\vec{n}\begin{pmatrix}2 \\ 1 \\ 2\end{pmatrix}$ est  un vecteur normal au plan (EFG), donc le plan (EFG) a une équation de la forme $2x+y+2z+d=0$ où $d$ est un réel à déterminer.

F est un point de (EFG) donc
$2x_{\text F} + y_{\text F} +2z_{\text F}+d=0$
c'est-à-dire\\
$2\times (-1) + 2 +2\times 1+d=0$ ce qui donne $d=-2$.

Le plan (EFG) a donc pour équation $2x+y+2z-2=0$.

		\item La droite $(d)$ passant par le point $\mathrm{D}$ et orthogonale au plan (EFG)  a pour vecteur directeur tout vecteur normal au plan (EFG) donc en particulier le vecteur $\vect{n}$. Elle a donc pour représentation paramétrique;

$\left \{
\begin{array}{l !{=} l l}
x & x_{\text D} + x_{\vec n}t\\
y & y_{\text D} +y_{\vec n}t & t\in\R\\
z & z_{\text D} + z_{\vec n}t
\end{array}
\right .	$	
c'est-à-dire
$\left \{
\begin{array}{l !{=} l l}
x & 3 + 2t\\
y & 1 +t & t\in\R\\
z & 5 + 2t
\end{array}
\right.$

		\item On note K le projeté orthogonal du point D sur le plan (EFG). Le point K appartient donc à la droite $(d)$ et au plan (EFG) donc ses coordonnées sont solutions du système:
$\left \{
\begin{array}{r !{=}  l}
x & 3 + 2t\\
y & 1 +t \\
z & 5 + 2t\\
2x+y+2z-2 & 0
\end{array}
\right.$

On a donc $t$ tel que $2\left (3+2t\right ) + \left (1+t\right ) + 2\left (5+2t\right ) -2=0$, 
soit\\
$6+4t+1+t+10+4t-2 = 0$
ou encore
$9t+15 = 0$; ce qui donne $t=-\dfrac{15}{9}=-\dfrac{5}{3}$.
		
On a donc:
$x_{\text K} =3+2t=3+2\left (-\dfrac{5}{3}\right ) = -\dfrac{1}{3}$,
$y_{\text K} = 1+t=1+\left (-\dfrac{5}{3}\right )=-\dfrac{2}{3}$ et
$z_{\text K} = 5+2t = 5+2\left (-\dfrac{5}{3}\right )=\dfrac{5}{3}$

Donc le point K a pour coordonnées $\left ( -\dfrac{1}{3}~,~ -\dfrac{2}{3} ~,~\dfrac{5}{3} \right )$.


	\end{enumerate}
\item 
	\begin{enumerate}
		\item 
$\begin{gathered}\text{DK}^2
= \left (-\dfrac{1}{3}-3 \right )^2 +\left (-\dfrac{2}{3}-1 \right )^2 +\left (\dfrac{5}{3}-5 \right )^2
= \left (-\dfrac{10}{3} \right )^2 +\left (-\dfrac{5}{3} \right )^2 +\left (-\dfrac{10}{3} \right )^2\\[7pt]
\phantom{\text{DK}^2}
=  \dfrac{100}{9} + \dfrac{25}{9}+\dfrac{100}{9} = \dfrac{225}{9}=25\end{gathered}$

Donc la distance DK est égale à 5~cm.
		
		\item Le volume du tétraèdre DEFG est donné par la formule:
$ \dfrac{\text{aire de la base}\times\text{hauteur}}{3}$.

Une base du tétraèdre est le triangle EFG d'aire 12,
et la hauteur relative à cette base est [DK] de longueur 5.				

Le volume du tétraèdre DEFG est donc, en cm$^3$:
$\dfrac{12\times 5}{3}=20$.
	\end{enumerate}
\end{enumerate}

\bigskip

\end{document}
